\section{客观违法阻却事由 - 正当防卫}
\begin{enumerate}
  \item 正当防卫。制止不法侵害。不能狭隘理解为人拿刀砍我，我反击！！
  \item 客观要件 \yjt 客观违法阻却事由 \yjt | 主观要件。所以先要有违法性，才能阻却掉。命题陷阱。
        \al{甲乙在农贸市场内\xhx{嬉戏打闹}}。
  \item 成立条件。5 个。起因、时间、意思、对象、限度。
        \begin{enumerate}
          \item 起因条件。现实的不法侵害。
                \begin{enumerate}
                  \item 主要针对个人法益，社会/国家法益由公权力处理。\al{甲看见乙正在\xhx{翻越国边境}，开枪重伤乙}。
                  \item 不法侵害具有客观性。\al{精神病人/未成年打人 \yjt 如果可以先采取避免措施}。
                  \item 现实性。不能是假想。
                        \begin{kk}{假想防卫}
                          不存在现实的不法侵害，客观 \cuo，主观 \dui \yjt 不成立正当防卫，有过失定过失，没有定意外事件。
                          \al{甲入睡时发现窗外有人影，以为是小偷要行窃。直接扔一块砖。原来是女儿男朋友幽会}。
                        \end{kk}
                \end{enumerate}
          \item 时间条件。适时。正在进行的不法侵害。
                \begin{enumerate}
                  \item 开始。设立防卫装置。不能危害公共安全 \yjt 防卫过当。\al{看家护院的藏獒}。
                  \item 结束。
                        \begin{enumerate}
                          \item 2 项要素。行为举止要结束 + 不法侵害危险性要解除（反复的可能性）。一般人行为时判断。
                          \item 财产犯罪里取得了财物，但被害人能当场挽回的，视为没结束。\al{追上小偷把小偷打趴下了}。
                        \end{enumerate}
                \end{enumerate}
          \item 防卫意思。是一个主观条件。
                \begin{enumerate}
                  \item 不要求具有防卫意图/意志/动机。防卫人制止不法侵害是出于何种目的。只考了一次。
                        \begin{enumerate}
                          \item \al{甲看到仇人乙在强制猥亵妇女，\xhx{出于报复}心理攻击乙，将乙打倒}。
                          \item \al{甲看到仇人乙在强制猥亵自己的女友，一方面\xhx{报复心理}，一方面\xhx{救助心理}，将乙打倒}。
                        \end{enumerate}
                  \item 防卫认识/意识。
                        \begin{kk}{偶然防卫。观点展示}
                          制止不法侵害上没做到主客观一致，客观 \dui，主观 \cuo。
                          \begin{enumerate}
                            \item 故意型。\al{乙杀丙，甲杀乙，结果救了丙}。
                                  \begin{itemize}
                                    \item 不要说 \yjt 结果无价值 \yjt 甲正当防卫。
                                    \item 必要说 \yjt 行为无价值 \yjt 故意杀人未遂。既遂要求坏结果，现在是好结果，只能定未遂。
                                  \end{itemize}
                            \item 过失型。\al{乙杀丙，甲擦枪不慎走火打到乙}。
                                  \begin{itemize}
                                    \item 不要说。正当防卫 \yjt 无罪。
                                    \item 必要说。不成立正当防卫 \yjt 不成立过失致人死亡（因为是一个好结果） \yjt 无罪。
                                  \end{itemize}
                          \end{enumerate}
                          \tcblower
                          这里的行为无价值应理解成不评价行为，只评价结果好坏。\al{枪杀犯人 \yjt 好结果}。
                        \end{kk}
                  \item 防卫挑拨。恶意利用正当防卫。\al{相互斗殴，狗咬狗 \yjt 谁先动手谁丧失正当防卫资格}。
                \end{enumerate}
          \item 对象条件。对不法侵害本人施加防卫手段。
                \begin{enumerate}
                  \item 共同犯罪。对正犯可以，帮助犯看是否有攻击性。
                \end{enumerate}
          \item 限度条件。防卫手段过当问题。客观必须发生过当结果，主观至少要有过失。
                \begin{enumerate}
                  \item 必要性优先说。应作为第一标准。只要是制止不法性的必要手段，就不过当。\al{多次强奸寡妇被毒死}。
                  \item 相当性优先说。比例原则。
                        \begin{enumerate}
                          \item \al{甲入户盗窃，卡在窗外，主人乙将甲砍成重伤 \yjt 财产法益 \xyh 人身法益}。
                          \item \al{甲得知乙今晚会再次从阳台入室盗窃，设立防卫装置，乙翻入时，防卫装置启动，乙人头落地}。
                        \end{enumerate}
                  \item vs 事后防卫。防卫过当前提是正当防卫成立。
                        \al{甲伤害乙，乙奋起反抗，导致甲头撞桌角，\xhx{倒地}昏迷，嘴角抽搐，乙\xhx{见状抡起小板凳}砸甲的头}。
                \end{enumerate}
        \end{enumerate}
  \item 特殊正当防卫。\al{行凶、杀人、抢劫、强奸、绑架、其他\xhx{严重危及}人身安全的\xhx{暴力犯罪}（\xhx{犯罪行为而非罪名}）}。
  \item 结合刑诉扭送行为。\al{甲持刀抢劫超市，乙夺刀扔掉，甲未取得财物见状\xhx{逃走}，\xhx{乙追上}扑倒甲摔成重伤}。
  \item 做正当防卫、紧急避险的题剥离各种修饰语、身份等主观东西，只抽取事实（抹掉故意、过失）。
        \al{甲看到\scx{仇人}乙在强制猥亵\scx{自己的}女\scx{友，一方面出于报复心理，一方面出于救助心理}，将乙打倒}。
\end{enumerate}
